\chapter*{Ход выполнения практики}
\addcontentsline{toc}{chapter}{Ход выполнения практики}

\section*{Задания и отметки о выполнении}

\task{Описание математической модели динамики роста опухоли с учетом иммунотерапии на
основе системы дифференциальных уравнений.}

\task{Приведение численных методов, используемых для решения данной системы уравнений.}

\task{Проектирование и реализация программного обеспечения с графическим
пользовательским интерфейсом для моделирования динамики роста опухолей.}

\task{Проведение серии вычислительных экспериментов для установления влияния
различных входных параметров модели на динамику роста опухоли.}

\task{Подготовка сравнительного анализа полученных результатов численных
экспериментов с данными клинических наблюдений.}

\task{Подготовка отчёта о проделанной работе в системе LaTeX.}

\task{Создание презентации в соответствии с данным отчетом и устного доклада
для публичной защиты.}


% Для вкр не нужен

% \clearpage

% \section*{Отзыв о работе обучающегося}

% В период прохождения практики Карманчиков В.А. провел анализ предметной
% области, связанной с компьютерным моделированием динамики опухолевых процессов
% при иммунотерапии. В ходе работы были изучены подходы к математическому
% описанию взаимодействия опухолевых клеток, иммунных эффекторных клеток и
% цитокинов на основе систем обыкновенных дифференциальных уравнений, а также
% рассмотрены численные методы решения таких систем.

% Владислав Андреевич разработал концепцию программного обеспечения для
% моделирования динамики опухолевых процессов, определил структуру математической
% модели, начальные параметры системы и основные сценарии терапевтического
% воздействия. Им была выполнена программная реализация вычислительного ядра на
% языке Python с использованием библиотек SciPy, NumPy, Matplotlib и PyQt5, а
% также разработан графический пользовательский интерфейс для проведения
% численных экспериментов и визуализации результатов.

% В рамках работы были проведены вычислительные эксперименты, включающие
% валидацию модели на экспериментальных данных, параметрический анализ
% чувствительности системы, а также моделирование различных терапевтических
% стратегий, включая введение цитокинов, адоптивный перенос клеток и
% комбинированные схемы терапии.

% В целом поставленные перед Карманчиковым В.А. задачи выполнены, а
% предусмотренные учебным планом компетенции освоены на достаточном уровне.
% Считаю, что Карманчиков Владислав Андреевич заслуживает оценки <<удовлетворительно>>.

\progressapprov
